% `tex.web` distinguishes primitives by `(cmd, chr)`. This engine has one
% variant per *command*, and four families were collapsed onto one variant
% each — so which member a control sequence was had to be recovered from the
% way it was spelled:
%
%     match t.as_cs().map(resolve).as_deref() {
%         Some("displaystyle") => Style::Display,
%         …
%         _ => Style::ScriptScript,   // the fallback
%     }
%
% A `\let` alias spells something else, so it fell to the `_` arm. `\let\ts=
% \textstyle` selected \scriptscriptstyle — silently, and on the page — and
% `\let\nl=\nolimits` meant \displaylimits. expl3 renames every primitive it
% uses into its own namespace and then undefines the original name, so this is
% not a corner case; it is what happens to a document that loads expl3.
%
% §1171's four styles and §1156's three limit switches are the two that reach
% the page. §1071's four movers are the same shape with a milder consequence:
% the *dispatch* was already right, because `\raise` and `\moveleft` differ
% only in a sign this engine had, but `\meaning` guessed from the current mode
% — which is what §367's trace should do, since the modes never overlap, and
% is arbitrary for `\meaning`, which can be asked anywhere.
%
% Widths are the test for the styles because a script style picks a smaller
% font, and the same `x` set in three sizes has three widths.
\input plain
\font\rm=cmr10 \font\rs=cmr7 \font\rss=cmr5
\textfont0=\rm \scriptfont0=\rs \scriptscriptfont0=\rss
\font\mi=cmmi10 \font\mis=cmmi7 \font\miss=cmmi5
\textfont1=\mi \scriptfont1=\mis \scriptscriptfont1=\miss
\font\sy=cmsy10 \font\sys=cmsy7 \font\syss=cmsy5
\textfont2=\sy \scriptfont2=\sys \scriptscriptfont2=\syss
\rm
\let\ds=\displaystyle \let\ts=\textstyle
\let\ss=\scriptstyle  \let\sss=\scriptscriptstyle
\let\li=\limits \let\nl=\nolimits \let\dl=\displaylimits
\let\ml=\moveleft \let\mr=\moveright \let\ra=\raise \let\lo=\lower

% Each style, written out and through an alias. The two must agree.
\setbox0\hbox{$\displaystyle x$}\setbox1\hbox{$\ds x$}
\setbox2\hbox{$\textstyle x$}\setbox3\hbox{$\ts x$}
\setbox4\hbox{$\scriptstyle x$}\setbox5\hbox{$\ss x$}
\setbox6\hbox{$\scriptscriptstyle x$}\setbox7\hbox{$\sss x$}
\message{<ds=\the\wd0|\the\wd1>}
\message{<ts=\the\wd2|\the\wd3>}
\message{<ss=\the\wd4|\the\wd5>}
\message{<sss=\the\wd6|\the\wd7>}

% ...and the three limit switches, where a big operator's script goes above it
% or beside it.
\setbox8\hbox{$\displaystyle\mathop{\rm A}\limits^2$}
\setbox9\hbox{$\displaystyle\mathop{\rm A}\li^2$}
\setbox10\hbox{$\displaystyle\mathop{\rm A}\nolimits^2$}
\setbox11\hbox{$\displaystyle\mathop{\rm A}\nl^2$}
\setbox12\hbox{$\displaystyle\mathop{\rm A}\displaylimits^2$}
\setbox13\hbox{$\displaystyle\mathop{\rm A}\dl^2$}
\message{<li=\the\wd8|\the\wd9>}
\message{<nl=\the\wd10|\the\wd11>}
\message{<dl=\the\wd12|\the\wd13>}

% The movers still shift the way they always did, alias or not.
\setbox14\vbox{\moveleft 5pt\hbox{x}}
\setbox15\vbox{\ml 5pt\hbox{x}}
\setbox16\hbox{\raise 5pt\hbox{x}}
\setbox17\hbox{\ra 5pt\hbox{x}}
\message{<mv=\the\ht14|\the\ht15|\the\ht16|\the\ht17>}

% §298 names each of them for itself, through an alias too.
\message{<s=\meaning\displaystyle|\meaning\textstyle|\meaning\scriptstyle|\meaning\scriptscriptstyle>}
\message{<sa=\meaning\ds|\meaning\ts|\meaning\ss|\meaning\sss>}
\message{<l=\meaning\limits|\meaning\nolimits|\meaning\displaylimits>}
\message{<la=\meaning\li|\meaning\nl|\meaning\dl>}
\message{<v=\meaning\moveleft|\meaning\moveright|\meaning\raise|\meaning\lower>}
\message{<va=\meaning\ml|\meaning\mr|\meaning\ra|\meaning\lo>}
\bye
